\documentclass[10pt]{article}
\usepackage[margin=0.9in]{geometry}
\usepackage{amsmath,amssymb,booktabs,hyperref}
\title{Conditional Predictive Evidence for Diagnostic Reuse:\\An Initial Controlled Study}
\author{Working research manuscript}
\date{18 October 2026}
\begin{document}
\maketitle
\begin{abstract}
A useful intervention need not validate the explanation that motivated it.
We implement a data-intervention learner with separate updates for observed recipe value and diagnostic predictive evidence.
The latter compares local-response predictions conditional on total improvement and affects reuse for unqueried candidates.
We evaluate the implementation retrospectively using 96 existing seed-level parent--child comparisons, aggregated into 32 decisions across five tasks from two source families.
Seventy fixed-expert policy trajectories and thirty generated-feature trajectories provide no decision advantage for conditional over joint evidence at the registered budgets.
Simple strong response models outperform the diagnostic mixtures.
Four source-isolated Sol-high generation calls yield four ordinary feature combinations and two abstentions from mismatch-driven generation.
These results establish an executable, auditable mechanism, not a successful new method or independent generalization.
\end{abstract}

\section{Question and scope}
We study data interventions under fixed forecasting backbones and source-specific frozen training protocols.
The decision problem is whether a diagnostic explanation should influence future unqueried interventions, while retaining useful recipes independently of their explanations.
Intervention-driven verification is already central to DoVer~\cite{dover}, and language-guided hypothesis and slice discovery appears in LADDER~\cite{ladder}.
Neither separate evaluation of repair and explanation nor a conditional log score is claimed as a novel mathematical tool.
The empirical contribution sought is additional decision value from controlling diagnostic reuse; the present study does not establish that contribution.

\section{An explicit separation of updates}
Let $L_P>0$ and $L_C$ be parent and child global losses under fixed evaluation weights.
A target region $S$ has mass $p\in(0,1)$ and complement $\bar S$.
Region losses are conditional means, while all response coordinates share the global parent denominator:
\[
 d_S=\frac{L_{P,S}-L_{C,S}}{L_P},\quad
 d_{\bar S}=\frac{L_{P,\bar S}-L_{C,\bar S}}{L_P},\quad
 Y=p\,d_S+(1-p)d_{\bar S},\quad Z=d_S-d_{\bar S}.
\]
Thus $d_S=Y+(1-p)Z$ and $d_{\bar S}=Y-pZ$.
These are two response coordinates, not three independent pieces of evidence.

Before observing a candidate, freeze reference and diagnostic predictors
$p_0(Y\mid x),p_0(Z\mid Y,x)$ and $p_h(Y\mid x),p_h(Z\mid Y,x)$.
After observing $(y,z)$, score
\[
 E_h^{\mathrm{cond}}=\log p_h(z\mid y,x)-\log p_0(z\mid y,x).
\]
Actual $y$ is a scoring input only; unqueried candidate selection never receives it.
With the same factorization,
\[
 E_h^{\mathrm{joint}}=E_h^Y+E_h^{\mathrm{cond}}.
\]
If both predictors share their $Y$ marginal, conditional and joint scores coincide exactly.
A meaningful joint-score ablation therefore permits different $Y$ marginals and reports the marginal contribution.

We accumulate $\rho_h\leftarrow\rho_h+E_h$ after scoring the frozen prediction.
Unit prior mass is assigned to the reference and unit total mass to all diagnostic experts:
\[
 w_0,w_1,\ldots,w_H\ \propto\
 1,\ \exp(\rho_1)/H,\ldots,\exp(\rho_H)/H.
\]
The acquisition score is $w_0\hat Y_0(x)+\sum_h w_h\hat Y_h(x)$.
All recipe regressors learn from actual revealed $Y$.
The best observed eligible recipe is selected using measured benefit, independently of $\rho$.
An event can contribute evidence once; a revised diagnostic receives a new version without inherited evidence.
These scores are predictive comparisons, not root-cause probabilities, Bayes factors, calibrated tests, or causal identification.

\section{Frozen development experiment}
Hydraulic contributes three tasks with four actions and two implementations, each with three seeds.
ETT contributes two stations with two actions and two implementations, each with three seeds.
Sampling and corresponding uniform-sampling loss weighting are distinct implementations.
The 96 seed responses form 32 candidate decisions by averaging per-seed normalized responses.
All results were already revealed historically; these are retrospective development experiments.
ETT stations belong to one family and Hydraulic tasks to another.

The target region is the latter half of the forecast horizon, across every evaluation window and channel, so $p=1/2$.
We use original validation-selected checkpoints and never reselect a checkpoint on candidate evaluation results.
A shared set of 22 parent-state and action primitives is supplied to every reference.
Four manual experts add fixed interactions.
Task-balanced ridge regression uses coefficient 1; leave-task-out residuals estimate the Gaussian scale, floored at 0.01 in response-ratio units.
We compare reference, all fixed features, all features with quadratic terms, and expert mixtures with uniform, marginal, joint, or conditional weights.

Each policy registers predictions, queries one three-seed candidate, scores its frozen prediction, then refits using only revealed records.
At 2 and 4 queries, the best observed candidate must have mean $Y\geq0.01$ and NMAE gain $\geq-0.01$, otherwise the parent is retained.
A query represents three logical full training calls, plus a shared three-seed parent.
Eight-query results are supplementary and cover Hydraulic only.
Both task holdout and family holdout are development stress tests, not final unseen-source confirmation.

\section{Results}
\begin{table}[h]
\centering
\caption{Family-held-out retrospective mean improvement (\%) over five tasks.}
\begin{tabular}{lrr}
\toprule
Policy & 2 queries & 4 queries\\
\midrule
Primitive-feature reference & 6.470222 & 7.798965\\
All fixed features & 7.332344 & 12.554207\\
All fixed features, quadratic & 6.984083 & 12.330422\\
Manual experts, uniform & 6.918507 & 7.798965\\
Manual experts, joint & 6.918507 & 7.798965\\
Manual experts, conditional & 6.918507 & 7.798965\\
Ordinary generated features, conditional & 6.470222 & 7.798965\\
Mismatch generator, conditional & 6.470222 & 7.798965\\
\bottomrule
\end{tabular}
\end{table}
The finite-pool oracle mean is 12.554207\%.
Conditional, joint, marginal, and uniform manual mixtures give the same recommendations at both registered budgets.
Conditional and joint scores do differ numerically; one family-held-out task changes its query order without changing selected benefit.
Task-held-out four-query means are 7.798965\% for the conditional mixture and 12.554207\% for the quadratic baseline.
Thus numerical evidence separation does not establish decision value.

Two source-only generation conditions, ordinary prediction-error feature engineering and response-mismatch diagnostic generation, each receive one call per source with identical facts and expression budgets.
The requested model is \texttt{gpt-6.1-sol} with high reasoning.
Ordinary generation returns two features in each direction; mismatch generation abstains in both.
Generated expressions only combine existing measurements and are not discoveries of additional raw information.
The mismatch generator identifies the lack of conditional-$Z$ residual evidence and the risk of treating task profiles as reusable mechanisms.

Ordinary combinations reduce static Hydraulic-to-ETT MAE from 32.075868 to 11.796932 percentage points, without improving the reported selection budgets.
ETT-to-Hydraulic prediction remains severely miscalibrated: MAE changes from 17414.065852 to 11006.237499 percentage points, while nominal 95\% coverage changes from 12.5\% to 8.33\%.
These values are retained without clipping.
Predictor extrapolation and measurement insufficiency cannot be conflated.

\section{Audit, cost, and limitations}
Nine core checks and an independent 70-trajectory audit pass.
The latter checks 448 steps, 2,944 regression component snapshots, 1,792 candidate predictions, 1,024 score triples, and 182 budget nodes.
Snapshot training identities contain only initially permitted and already revealed records.
The built-in hidden-label test redundantly invokes identical inputs and does not independently establish intervention invariance; the independent identity audit supplies the actual recorded access check.
A separate generated-policy audit checks 192 additional steps and 78 budget nodes, with 13,391 assertions passing. Thirty-two static generated-combination predictions lack saved fit snapshots; only their error and coverage arithmetic is independently checked.

No new full forecast training is performed; the historical completed count remains 464.
Historical candidate costs are known for 31 of 32 decisions and 146 of 182 manual budget nodes, so complete wall-clock equal-budget superiority cannot be claimed.
Four LLM calls consume 108,492 input and 10,599 output tokens, with 700.481619 summed call-seconds under parallel execution.
Transport retries and an initial validator error are retained; supplementary validation changes event classification, not generated outputs.
No test or sealed audit partition is evaluated.

The study has only five correlated tasks and two source families.
Positive means do not imply seed-stable benefit.
We have not evaluated a full strong-prompt adaptive agent, genuinely new measurement programs, a frozen-prefix feedback intervention, or independent-source confirmation.
The current evidence does not justify a claim of method superiority.
Further development must address predictor applicability and conditional response mismatch before expanding full-training experiments.

\begin{thebibliography}{9}
\bibitem{dover}
M. Ma et al.
\emph{DoVer: Intervention-Driven Auto Debugging for LLM Multi-Agent Systems}.
arXiv:2512.06749, 2025.
\url{https://arxiv.org/abs/2512.06749}.
\bibitem{ladder}
S. Ghosh et al.
\emph{LADDER: Language-Driven Slice Discovery and Error Rectification in Vision Classifiers}.
Findings of ACL, 2025; arXiv:2408.07832v13.
\url{https://arxiv.org/abs/2408.07832}.
\end{thebibliography}
\end{document}
